% Redakcyjny spis najczęstszych symboli (bez powtarzania definicji z rozdziałów).
\chapter*{Spis oznaczeń}
\markboth{SPIS OZNACZEŃ}{SPIS OZNACZEŃ}
\phantomsection\addcontentsline{toc}{chapter}{Spis oznaczeń}
\noindent Symbole mogą przyjmować dodatkowe założenia w poszczególnych twierdzeniach.
Pierwszeństwo mają zawsze definicje podane we wskazanym rozdziale.

\section*{Zbiory i mierzalność}
\begin{longtable}{@{}p{0.285\textwidth}p{0.555\textwidth}p{0.095\textwidth}@{}}
$X,\,\mathcal P(X)$ & Przestrzeń bazowa i zbiór jej wszystkich podzbiorów. & rozdz.~\ref{ch:przestrzenie-mierzalne} \\
$\mathcal A,\,\mathcal F$ & Rodziny zbiorów; w szczególności $\sigma$-algebry. & rozdz.~\ref{ch:przestrzenie-mierzalne} \\
$\sigma(\mathcal G)$ & Najmniejsza $\sigma$-algebra zawierająca rodzinę $\mathcal G$. & rozdz.~\ref{ch:przestrzenie-mierzalne} \\
$\mathcal B(X)$ & $\sigma$-algebra borelowska przestrzeni topologicznej $X$. & rozdz.~\ref{ch:przestrzenie-mierzalne} \\
$\chi_A$ & Funkcja charakterystyczna zbioru $A$; jedyna stosowana w książce notacja tego pojęcia. & rozdz.~\ref{ch:funkcje-mierzalne} \\
$\mu^*$ & Miara zewnętrzna (niekoniecznie zdefiniowana na $\sigma$-algebrze). & rozdz.~\ref{ch:caratheodory} \\
$\lambda_d$ & $d$-wymiarowa miara Lebesgue'a na $\mathbb R^d$. & rozdz.~\ref{ch:lebesgue} \\
$\mathrm d\mu(x)$ & Element notacji całkowania wskazujący miarę i zmienną całkowania. & rozdz.~\ref{ch:calkowanie-lebesgue} \\
\end{longtable}

\section*{Zbieżność i całkowanie}
\begin{longtable}{@{}p{0.285\textwidth}p{0.555\textwidth}p{0.095\textwidth}@{}}
$\mu\otimes\nu$ & Miara produktowa na mierzalnej przestrzeni produktowej. & rozdz.~\ref{ch:produkty-miary} \\
$L^p(\mu)$ & Przestrzeń klas funkcji $p$-całkowalnych; dla $p=\infty$ funkcji istotnie ograniczonych. & rozdz.~\ref{ch:przestrzenie-lp} \\
$\|f\|_{L^p}$ & Norma w przestrzeni $L^p(\mu)$. & rozdz.~\ref{ch:przestrzenie-lp} \\
$C_c(X)$ & Przestrzeń ciągłych funkcji o zwartym nośniku. & rozdz.~\ref{ch:radon-measures-20} \\
$C_0(X)$ & Przestrzeń ciągłych funkcji zanikających w nieskończoności. & rozdz.~\ref{ch:riesz-markov-21} \\
$F_*\mu$ & Miara obrazowa: $(F_*\mu)(A)=\mu(F^{-1}(A))$. & rozdz.~\ref{ch:rn-applications-19} \\
\end{longtable}

\section*{Miary i ich rozkłady}
\begin{longtable}{@{}p{0.285\textwidth}p{0.555\textwidth}p{0.095\textwidth}@{}}
$\nu\ll\mu$ & Bezwzględna ciągłość miary $\nu$ względem $\mu$. & rozdz.~\ref{ch:ac-singular-17} \\
$\nu\perp\mu$ & Wzajemna osobliwość miar. & rozdz.~\ref{ch:miary-ze-znakiem} \\
$\nu^+,\,\nu^-$ & Dodatnia i ujemna część rozkładu Jordana miary ze znakiem. & rozdz.~\ref{ch:miary-ze-znakiem} \\
$|\nu|$ & Wariacja całkowita miary ze znakiem, zespolonej lub wektorowej; zależnie od kontekstu. & rozdz.~\ref{ch:miary-ze-znakiem} \\
$\frac{\mathrm d\nu}{\mathrm d\mu}$ & Pochodna Radona--Nikodyma, a nie symbol całkowania. & rozdz.~\ref{ch:radon-nikodym-18} \\
$\operatorname{supp}\mu$ & Nośnik miary, z uwzględnieniem odpowiednich założeń topologicznych. & rozdz.~\ref{ch:radon-measures-20} \\
\end{longtable}

\section*{Miary geometryczne}
\begin{longtable}{@{}p{0.285\textwidth}p{0.555\textwidth}p{0.095\textwidth}@{}}
$\mathcal H^s_\delta,\,\mathcal H^s$ & Miary Hausdorffa na skali $\delta$ oraz po przejściu do granicy; z ustaloną normalizacją. & rozdz.~\ref{ch:hausdorff-26} \\
$\dim_{\mathrm H}(E)$ & Wymiar Hausdorffa zbioru $E$. & rozdz.~\ref{ch:hausdorff-26} \\
$\mu_g$ & Miara objętości związana z metryką riemannowską $g$. & rozdz.~\ref{ch:miary-rozmaitosci-27} \\
$J_k(A)$ & $k$-wymiarowy jakobian odwzorowania liniowego $A$. & rozdz.~\ref{ch:area-coarea-28} \\
$\operatorname{Jac}_m(B)$ & $m$-wymiarowy jakobian koobszaru odwzorowania liniowego $B$. & rozdz.~\ref{ch:area-coarea-28} \\
$BV(\Omega)$ & Przestrzeń funkcji o ograniczonej wariacji w $\Omega$. & rozdz.~\ref{ch:geometric-measure-29} \\
$BV_{\mathrm{loc}}(\Omega)$ & Przestrzeń funkcji o wariacji lokalnie skończonej. & rozdz.~\ref{ch:geometric-measure-29} \\
$Du,\,|Du|$ & Miara pochodnej dystrybucyjnej i jej wariacja. & rozdz.~\ref{ch:geometric-measure-29} \\
$P(E;\Omega)$ & Obwód zbioru $E$ względem otwartego zbioru $\Omega$. & rozdz.~\ref{ch:geometric-measure-29} \\
$\partial^*E,\,\nu_E$ & Brzeg zredukowany i zewnętrzna normalna miarowo-teoretyczna. & rozdz.~\ref{ch:geometric-measure-29} \\
\end{longtable}

\section*{Miary wektorowe}
\begin{longtable}{@{}p{0.285\textwidth}p{0.555\textwidth}p{0.095\textwidth}@{}}
$\nu:\mathcal A\to X$ & Miara o wartościach w przestrzeni Banacha $X$. & rozdz.~\ref{ch:vector-measures-30} \\
$\|\nu\|_{\mathrm{sv}}$ & Półwariacja miary wektorowej. & rozdz.~\ref{ch:vector-measures-30} \\
$\int_A g(x)\,\mathrm d\mu(x)$ & Całka funkcji wektorowej w sensie Bochnera, gdy $g$ spełnia odpowiednie założenia. & rozdz.~\ref{ch:vector-measures-30} \\
$\mathrm{RNP}$ & Własność Radona--Nikodyma przestrzeni Banacha. & rozdz.~\ref{ch:vector-measures-30} \\
\end{longtable}

\noindent\emph{Uwaga.} Notacja $\chi_A$ jest stosowana konsekwentnie
w całym skrypcie; miary geometryczne stosują wspólną kalibrację
ustaloną w rozdziale~\ref{ch:hausdorff-26}.


\section*{Dżety i przedłużanie danych}
\begin{longtable}{@{}p{0.285\textwidth}p{0.555\textwidth}p{0.095\textwidth}@{}}
$j_p^kf$ & Dżet $k$-tego rzędu funkcji w punkcie $p$. & rozdz.~\ref{ch:dzety} \\
$J^kE$ & Wiązka dżetów $k$-tego rzędu przekrojów $E\to M$. & rozdz.~\ref{ch:dzety} \\
$D_{\mathrm{ap}}f$ & Pochodna przybliżona, określona w odpowiednich punktach. & rozdz.~\ref{ch:dzety} \\
\end{longtable}
